\usepackage{fontspec}
\setmainfont{texgyretermes}[
  Extension      = .otf,
  UprightFont    = *-regular,
  BoldFont       = *-bold,
  ItalicFont     = *-italic,
  BoldItalicFont = *-bolditalic
]
\setsansfont{texgyreheros}[
  Extension      = .otf,
  UprightFont    = *-regular,
  BoldFont       = *-bold,
  ItalicFont     = *-italic,
  BoldItalicFont = *-bolditalic,
  Scale          = 0.95
]
\setmonofont{texgyrecursor}[
  Extension      = .otf,
  UprightFont    = *-regular,
  BoldFont       = *-bold,
  ItalicFont     = *-italic,
  BoldItalicFont = *-bolditalic,
  Scale          = 0.88
]
\usepackage{polyglossia}
\setmainlanguage{vietnamese}
\setotherlanguage{english}

\usepackage{geometry}
\geometry{a4paper, top=2.5cm, bottom=2.5cm, left=3cm, right=2cm, headheight=15pt}
\usepackage{setspace}
\onehalfspacing
\usepackage{indentfirst}
\usepackage{fancyhdr}
\pagestyle{fancy}
\fancyhf{}
\fancyhead[L]{\small Phân tích \& Thiết kế thuật toán - Bài tập lớn}
\fancyhead[R]{\small\thepage}
\renewcommand{\headrulewidth}{0.4pt}

\usepackage{amsmath, amssymb, amsthm}
\usepackage{mathtools}

\usepackage{booktabs}
\usepackage{array}
\usepackage{multirow}
\usepackage{tabularx}
\usepackage{longtable}
\newcolumntype{L}{>{\raggedright\arraybackslash}X}
\newcolumntype{G}{>{\centering\arraybackslash}p{0.55cm}}
\usepackage{graphicx}
\usepackage{float}
\usepackage{needspace}
\usepackage{caption}

\usepackage{enumitem}
\usepackage{pifont}
\newcommand{\vong}[1]{\ding{\numexpr171+#1\relax}}

\usepackage{fancyvrb}
\usepackage{listings, xcolor}
\definecolor{codebg}{gray}{0.96}
\lstset{
  backgroundcolor = \color{codebg},
  basicstyle      = \ttfamily\small,
  keywordstyle    = \color{blue!70!black}\bfseries,
  commentstyle    = \color{green!45!black}\itshape,
  stringstyle     = \color{orange!80!black},
  numbers         = left,
  numberstyle     = \tiny\color{gray},
  numbersep       = 8pt,
  breaklines      = true,
  frame           = single,
  framesep        = 5pt,
  xleftmargin     = 18pt,
  showstringspaces= false,
  captionpos      = b,
  literate        = {->}{{$\rightarrow$}}1 {<=}{{$\leq$}}1 {>=}{{$\geq$}}1
}

\usepackage[most]{tcolorbox}
\newtcolorbox{promptbox}[1]{
  colback=gray!4, colframe=gray!55!black, boxrule=0.5pt,
  title={Prompt · #1}, fonttitle=\small\bfseries, fontupper=\small,
  sharp corners=downhill, breakable
}
\newtcolorbox{ghichu}[1][Lưu ý]{
  colback=yellow!8, colframe=orange!70!black, boxrule=0.6pt,
  title=#1, fonttitle=\bfseries
}

\usepackage{hyperref}
\hypersetup{
  colorlinks = true, linkcolor = blue!60!black,
  urlcolor   = cyan!70!black, pdfstartview = FitH
}

\newcommand{\T}[1]{T\!\left(#1\right)}
\newcommand{\Oh}[1]{\mathcal{O}\!\left(#1\right)}
\newcommand{\Om}[1]{\Omega\!\left(#1\right)}
\newcommand{\Th}[1]{\Theta\!\left(#1\right)}
\newcommand{\lgn}{\log_2 n}
\DeclarePairedDelimiter{\floor}{\lfloor}{\rfloor}
\DeclarePairedDelimiter{\ceil}{\lceil}{\rceil}
